% =====================================================================
% GA-SUM-02 · EXP-1 · Ingeniería de Requerimientos (ISR-401) · UTEQ
% Tema 7: Herramientas automatizadas de gestión de requisitos (Jira)
% Compilar: latexmk -pdf main.tex
% =====================================================================
\documentclass[aspectratio=169,11pt]{beamer}

% ---------- Idioma y codificación ----------
\usepackage[utf8]{inputenc}
\usepackage[T1]{fontenc}
\usepackage[spanish,es-noshorthands]{babel}
\usepackage{lmodern}

% ---------- Paquetes ----------
\usepackage{graphicx}
\usepackage{booktabs}
\usepackage{tabularx}
\usepackage{array}
\usepackage{xcolor}
\usepackage{tikz}
\usetikzlibrary{arrows.meta,positioning,calc}

\graphicspath{{figuras/}}

% ---------- Tema ----------
\usetheme{Madrid}
\usecolortheme{default}
\setbeamertemplate{navigation symbols}{}
\setbeamertemplate{bibliography item}[text]

\newcommand{\expone}[1]{\hfill{\scriptsize\textcolor{gray}{Expone: #1}}}

\tikzset{
  caja/.style={draw, rounded corners, fill=blue!8, minimum width=2.4cm,
               minimum height=0.9cm, align=center, font=\small},
  flecha/.style={-{Stealth[length=2.5mm]}, thick}
}

% ---------- Datos de la portada ----------
\title[Jira para la gestión de requisitos]{Herramientas automatizadas de apoyo a la IR: Jira}
\subtitle{GA-SUM-02 · EXP-1 · Aplicación al caso UteqFit-AI}
\author[Equipo AAM]{%
  Anderson Adonis Alcívar Vélez \\ Rai Eduardo Álvarez Rubio \\ Erick Ariel Monserrate Ochoa}
\institute[UTEQ]{Universidad Técnica Estatal de Quevedo\\
  Ingeniería de Requerimientos (ISR-401) · Cuarto semestre · Paralelo A}
\date{21 de noviembre de 2026}

\begin{document}

% =====================================================================
% 1. Portada
% =====================================================================
\begin{frame}
  \titlepage
\end{frame}

% --- Aquí van las diapositivas 2 a 5 (Rai Álvarez) ---

% --- Aquí van las diapositivas 6 a 9 (Anderson Alcívar) ---

% =====================================================================
% 10. Trazabilidad
% =====================================================================
\begin{frame}{Trazabilidad en Jira}
  \expone{Erick Monserrate}
  % Contenido pendiente
\end{frame}

% =====================================================================
% 11. Ventajas y limitaciones
% =====================================================================
\begin{frame}{Ventajas y limitaciones}
  \expone{Erick Monserrate}
  % Contenido pendiente
\end{frame}

% =====================================================================
% 12. Comparación con alternativas
% =====================================================================
\begin{frame}{Comparación con otras alternativas}
  \expone{Erick Monserrate}
  % Contenido pendiente
\end{frame}

% =====================================================================
% 13. Conclusiones
% =====================================================================
\begin{frame}{Conclusiones}
  \expone{Erick Monserrate}
  % Contenido pendiente
\end{frame}

% =====================================================================
% 14. Declaración de uso de IA
% =====================================================================
\begin{frame}{Declaración de uso de inteligencia artificial}
  \small
  \begin{tabularx}{\textwidth}{>{\bfseries}l X}
    \toprule
    Herramienta  & \\
    Versión      & \\
    Propósito    & \\
    Verificación & % COMPLETAR: cómo revisaron y reescribieron el contenido,
                   % y qué fuentes consultaron para comprobar cada afirmación.
                   \\
    \bottomrule
  \end{tabularx}
\end{frame}

% =====================================================================
% 15. Referencias (IEEE)
% =====================================================================
\begin{frame}[allowframebreaks]{Referencias}
  \scriptsize
  \bibliographystyle{IEEEtran}
  \bibliography{referencias}
\end{frame}

\end{document}
